\newpage
\section{Wnioski}

Projekt należy uznać za udany na poziomie inżynierskim i badawczym. Udało się zbudować kompletny pipeline do klasyfikacji obrazów \texttt{real} oraz \texttt{fake}, przygotować wariant globalny i twarzowy, wdrożyć narzędzia interpretowalności oraz przeprowadzić eksperymenty zarówno na wydzielonych zbiorach testowych, jak i na małym benchmarku OOD z nowymi generatorami.

Do najważniejszych osiągnięć należy zaliczyć:

\begin{itemize}
    \item poprawne uruchomienie treningu na GPU,
    \item uzyskanie wysokich wyników na bazowym zbiorze testowym,
    \item zbudowanie osobnego etapu twarzowego opartego o cropy twarzy,
    \item wdrożenie interpretowalności poprzez Grad-CAM,
    \item przygotowanie testu odporności na pogorszenie jakości,
    \item stworzenie lokalnego interfejsu demonstracyjnego,
    \item przygotowanie benchmarku OOD dla nowszych generatorów i przeprowadzenie analizy błędów.
\end{itemize}

\subsection{Co udało się osiągnąć}

Najważniejszym rezultatem projektu jest działający system, który obejmuje cały proces: przygotowanie danych, trening modeli, ewaluację, analizę odporności, interpretowalność oraz lokalny interfejs demonstracyjny. Model globalny osiągnął wysoką dokładność na głównym zbiorze testowym, a model twarzowy uzyskał bardzo dobre wyniki na danych portretowych zbliżonych do treningowych.

Udało się również pokazać, że sama dokładność na klasycznym zbiorze testowym nie jest wystarczająca do pełnej oceny systemu. Dlatego w projekcie wykonano dodatkowe testy odporności oraz benchmark OOD. Dzięki temu dokumentacja pokazuje nie tylko najlepsze wyniki, ale też ograniczenia modelu w trudniejszych warunkach.

\begin{table}[H]
\centering
\caption{Podsumowanie rezultatów i najważniejszych uwag}
\label{tab:cele-projektu}
\begin{tabular}{|L{4.2cm}|L{5.0cm}|L{4.2cm}|}
\hline
\textbf{Obszar} & \textbf{Co udało się osiągnąć} & \textbf{Ograniczenia lub uwagi} \\
\hline
Model globalny & Wysoka skuteczność na głównym zbiorze testowym: Accuracy 0.9339 dla baseline oraz 0.9268 dla wariantu robust. & Spadek skuteczności na benchmarku OOD do Accuracy 0.7125. \\
\hline
Odporność na jakość obrazu & Model zachował dobrą jakość dla lekkiego rozmycia i lekkiego downscale. & Silna kompresja JPEG oraz profil \texttt{mixed\_quality} wyraźnie obniżały wyniki. \\
\hline
Model twarzowy & Bardzo dobry wynik na wydzielonym teście portretowym: Accuracy 0.9763. & Wynik nie przenosił się automatycznie na nowe generatory spoza głównego rozkładu danych. \\
\hline
Adaptacja twarzowa & Poprawa po fine-tuningu na nowych cropach: Accuracy 0.8548 na teście adaptacyjnym. & Zbiór adaptacyjny był mały, więc wynik nie gwarantuje pełnej stabilności na dowolnych danych. \\
\hline
Interfejs demonstracyjny & Możliwość ręcznego testowania obrazów, porównania modeli i generowania Grad-CAM. & Interfejs spełnia założoną funkcję demonstracyjną i diagnostyczną. \\
\hline
\end{tabular}
\end{table}

Jednocześnie analiza praktyczna pokazała, że nawet bardzo dobry wynik ilościowy na danych zbliżonych do treningowych nie oznacza pełnej gotowości modelu do pracy na dowolnych danych. Szczególnie problematyczne okazały się:

\begin{itemize}
    \item nowe realistyczne portrety spoza głównego rozkładu treningowego,
    \item obrazy po silnej kompresji,
    \item najnowsze fotorealistyczne generatory obrazów,
    \item małe zbiory adaptacyjne, które nie pozwalają jeszcze na pełną stabilizację modelu,
    \item ograniczona dostępność dużych, publicznych i łatwo porównywalnych zbiorów obejmujących najnowsze zamknięte generatory obrazów.
\end{itemize}

\subsection{Co było problematyczne i dlaczego}

Największym problemem okazała się generalizacja, czyli przeniesienie skuteczności z wydzielonych zbiorów testowych zbliżonych do danych treningowych na nowe obrazy spoza głównego rozkładu. Przyczyną był przede wszystkim inny rozkład danych: nowe portrety i obrazy wygenerowane przez GPT / Gemini miały inną estetykę, inne artefakty oraz często wyższą jakość wizualną niż część danych treningowych.

Drugim problemem były fałszywe alarmy dla prawdziwych portretów. Model globalny bardzo dobrze wykrywał większość obrazów syntetycznych w benchmarku OOD, ale jednocześnie oznaczał wiele prawdziwych zdjęć jako \texttt{fake}. Możliwą przyczyną jest to, że model nauczył się traktować czyste, studyjne albo bardzo estetyczne portrety jako podejrzane, ponieważ podobne cechy często występują również w obrazach generowanych przez AI.

Trzecim problemem była wrażliwość na silne pogorszenie jakości obrazu. Kompresja JPEG, rozmycie i skalowanie mogą usuwać drobne artefakty, na których model opiera decyzję. W efekcie obraz po przetworzeniu przez komunikator albo serwis społecznościowy może być trudniejszy do klasyfikacji niż oryginalna próbka.

Osobnym ograniczeniem był model twarzowy. Analiza samej twarzy była przydatna, ale nie wystarczyła do pełnego rozwiązania problemu. Przyczyną jest to, że crop twarzy usuwa część kontekstu obrazu, a jednocześnie wymaga poprawnej detekcji twarzy. Dodatkowo mały zbiór adaptacyjny nie dawał wystarczająco dużo przykładów, aby model stabilnie nauczył się wszystkich wariantów nowych generatorów.

Istotnym ograniczeniem niezależnym od samej implementacji była także dostępność danych. Najnowsze modele generowania obrazów rozwijają się bardzo szybko, a publiczne zbiory referencyjne zwykle pojawiają się z opóźnieniem, są budowane dla innych celów albo nie obejmują dokładnie tych samych generatorów, stylów i typów obrazów, które są problematyczne w praktyce. Z tego powodu w projekcie nie dało się oprzeć oceny wyłącznie na dużym, gotowym i powszechnie uznanym benchmarku dla najnowszych modeli. Konieczne było przygotowanie własnego, mniejszego benchmarku OOD, który pozwolił sprawdzić kierunek problemu, ale nie zastępuje dużego zbioru referencyjnego.

Najważniejszy wniosek metodologiczny jest następujący: model globalny stanowi mocny baseline dla fotografii i znanych danych, natomiast generalizacja na najnowsze generatory GPT / Gemini pozostaje ograniczona. W benchmarku OOD problem nie polegał wyłącznie na niewykrywaniu obrazów syntetycznych, ale również na generowaniu zbyt wielu fałszywych alarmów dla nowych prawdziwych portretów.

Równie istotny jest wniosek dotyczący etapu twarzowego. Sam fakt zastosowania osobnego modelu twarzowego nie gwarantuje poprawy na nowych danych. Pierwotne checkpointy twarzowe nie generalizowały dobrze na nowe obrazy GPT / Gemini, a mała adaptacja na nowych cropach poprawiła wynik na teście adaptacyjnym. Nadal jednak ręczne testy pokazały trudne przypadki, w których nawet adaptowany model twarzowy klasyfikował syntetyczny portret jako \texttt{real}.

Oznacza to, że końcowy efekt pracy należy interpretować jako:

\begin{itemize}
    \item działający system badawczy z pełnym pipeline'em,
    \item mocny baseline dla starszych i znanych danych,
    \item narzędzie pozwalające badać ograniczenia detekcji na nowszych generatorach,
    \item punkt wyjścia do dalszych prac nad lepszą generalizacją modeli twarzowych.
\end{itemize}

W kolejnych pracach najbardziej uzasadnione byłyby:

\begin{enumerate}
    \item rozbudowanie zbioru OOD o większą liczbę obrazów z najnowszych generatorów, najlepiej po pojawieniu się większych publicznych benchmarków dla tych modeli,
    \item przygotowanie większego i bardziej zrównoważonego zbioru adaptacyjnego dla twarzy,
    \item sprawdzenie silniejszych metod detekcji twarzy i alternatywnych strategii cropowania,
    \item dalsza analiza modeli hybrydowych, które łączą sygnał globalny i twarzowy.
\end{enumerate}

Można zatem stwierdzić, że projekt osiągnął wyraźny i mierzalny progres: posiada sprawdzony etap bazowy, rozszerzenie twarzowe, benchmark OOD oraz udokumentowaną analizę ograniczeń. Z punktu widzenia pracy końcowej jest to wartościowy rezultat, ponieważ pokazuje nie tylko implementację, ale też realistyczną ocenę granic skuteczności obecnych metod detekcji.
